\subsection{Compression: De-duplication Impact on Prediction}
\label{sec:compression}

\paragraph{Task Description.}
The compression experiment ingests 200 samples (128\,B each,
25\,KB) and asks: does collision-based de-duplication degrade the
prediction machinery? After the substrate collapses colliding byte
patterns into shared attractors, $N = 50$ prefix queries are scored
by prediction-internal metrics—beam confidence, rejection rate, score
spread (hypothesis diversity), and origin diversity (mesh utilization).

\paragraph{Results.}
Compression ratio: 512.0\,bytes/entry (25\,KB $\to$ 50 entries).
Post-compression beam confidence remains stable across query depth
(top-score $$\mu = -0.007$, $\sigma = 0.019$, median $= -0.014$, range $[-0.030,\,0.031]$, $N = 50$$). Rejection rate: $\mu = 0.875$. Origin diversity:
$\mu = 1.0$ tries contributing.

Figure~\ref{fig:compression_chart} shows whether de-duplication
degrades beam confidence (top-left), increases mesh noise (top-right),
reduces hypothesis diversity (bottom-left), or narrows mesh
utilization (bottom-right).

\paragraph{Assessment.}
High compression (ratio 512.0) correlates with elevated rejection
($0.875$). De-duplication merges enough attractors that affinity routing
sends queries to tries with overlapping basins—the node-level beam
must prune more aggressively. This is the measurable cost of
compression on prediction quality.
